\metatitle{Rank functions and flats}
\metadescription{Joint rank, ordinary matroid ranks in each color,
conditional flats, common contractions, and best-completion valuations.}

\section[jointRank]{Rank functions and flats}

A color projection records which sets can receive one color. To retain
the compatibility between colors, we allow a separate set of positions
for each color and ask how many listed assignments a word can realize.
Fix a colored matroid $\mathcal F$ on $E$, with $n=|E|$ and quotas
$r_1,\dotsc,r_k$. We use $B(u,c)=\{e:u_e=c\}$ as in the
\href{intro.htm#coloredMatroids}{introduction}.

\begin{definition}[Joint rank]
For sets $S_1,\dotsc,S_k\subseteq E$, the \defin{joint rank} is
\[
 \defin{R(S_1,\dotsc,S_k)}\coloneqq
 \max_{u\in\mathcal F}\sum_{c=1}^k|B(u,c)\cap S_c|.
\]
The sets $S_c$ may overlap.
\end{definition}

With only the color-$c$ argument nonempty, this is the ordinary rank
of the projection matroid $M_c$. More generally, the joint rank
determines the entire family:
\[
 \begin{gathered}
 w\in\mathcal F\quad\Longleftrightarrow\\
 R(B(w,1),\dotsc,B(w,k))=n.
 \end{gathered}
\]
Indeed, the rank counts agreements with $w$, and $n$ agreements force
equality.

For a \href{independence.htm#independenceAndCircuits}{partial assignment}
$\tau$, put $T_c=\{e:(e,c)\in\tau\}$ and
$\rho(\tau)=R(T_1,\dotsc,T_k)$. Then $\rho(\tau)$ is the largest size
of an independent subset of $\tau$. In particular, $\tau$ is independent
if and only if $\rho(\tau)=|\tau|$, and a circuit $C$ has rank $|C|-1$.

\subsection[twoColorRank]{Two colors}

When colors $1$ and $0$ record a basis of $M$ and its complement,
the joint rank has the classical expression
\[
 \begin{aligned}
 R(A,B)={}&|A\cap B|+\operatorname{rk}_M(A\setminus B)\\
 &+\operatorname{rk}_{M^*}(B\setminus A).
 \end{aligned}
\]
Every position in $A\cap B$ contributes one. On the disjoint sets
$A\setminus B$ and $B\setminus A$, maximize the independent color-$1$
and color-$0$ prescriptions. They extend simultaneously by the
\href{independence.htm#twoColorIndependence}{two-color independence
criterion}.

\section[matroidRankSlices]{An ordinary rank in each color}

The key fact is that a unit bonus for one color can be optimized within
the maximizers of an integer score. Recall that
\href{operations.htm#scoreMaximizers}{score-maximizing words form a
colored matroid}.

\begin{proposition}[Adding a unit bonus]
Give every assignment an integer score $a_{e,d}$, and put
\[
 a(u)=\sum_{e\in E}a_{e,u_e},\qquad
 h=\max_{u\in\mathcal F}a(u).
\]
Let $M$ be the color-$c$ projection of the maximizing family. Then
for every $S\subseteq E$,
\[
 \begin{gathered}
 \max_{u\in\mathcal F}\bigl(a(u)+|B(u,c)\cap S|\bigr)\\
 =h+\operatorname{rk}_M(S).
 \end{gathered}
\]
\end{proposition}

\begin{proof}
Choose a word $u$ maximizing the combined score, with $a(u)$ largest
among ties. Suppose $a(u)<h$, and choose $v$ with $a(v)=h$.

The \href{exchange-graphs.htm#disjointCycles}{disjoint-cycle decomposition}
expresses $a(v)-a(u)>0$ as gains of individually feasible switches
from $u$. One cycle has integer gain at least one.

That cycle replaces at most one position of color $c$, losing at most
one unit of bonus. The combined score therefore does not decrease,
while $a(u)$ increases, contradicting the choice of $u$.
Thus an optimum lies among the words with $a(u)=h$, where the largest
bonus is $\operatorname{rk}_M(S)$.
\end{proof}

Now fix all arguments of $R$ except the color-$c$ argument, and write
\[
 R_c(S)=R(S_1,\dotsc,S_{c-1},S,S_{c+1},\dotsc,S_k).
\]

\begin{theorem}[Matroid rank slices]
There is an ordinary rank-$r_c$ matroid $M$ on $E$ such that
\[
 R_c(S)-R_c(\varnothing)=\operatorname{rk}_M(S).
\]
It is the color-$c$ projection of the words maximizing
$\sum_{d\ne c}|B(u,d)\cap S_d|$.
\end{theorem}

\begin{proof}
Apply the unit-bonus proposition to that integer score, whose maximum
is $R_c(\varnothing)$.
\end{proof}

In particular, $R$ is submodular in each argument separately. In a
fixed context, its ordinary matroid \defin{closure} is
\[
 \operatorname{cl}(S)=
 \{e\in E:R_c(S\cup\{e\})=R_c(S)\}.
\]
A \defin{flat} is a set $X$ with $\operatorname{cl}(X)=X$.
The flats form the usual geometric lattice $\mathcal L(M)$, with
intersection as meet and $\operatorname{cl}(X\cup Y)$ as join
\cite{Oxley2011}. Empty sets in the other color arguments recover
the flats of the original projection $M_c$.

\subsection[jointSubmodularity]{Joint submodularity fails}

Separate submodularity does not extend to assignments of several colors.
Take $\mathcal F=\{\mathtt{123},\mathtt{231}\}$ and the partial assignments
\[
 \begin{aligned}
 A&=\{(1,1),(3,1)\},\\
 B&=\{(2,2),(3,1)\}.
 \end{aligned}
\]
The two words differ on one three-color cycle, so this is a colored
matroid. We have
\[
 \begin{gathered}
 \rho(A)=\rho(B)=\rho(A\cap B)=1,\\
 \rho(A\cup B)=2.
 \end{gathered}
\]
Thus joint submodularity fails even when the combined prescriptions
have no conflicting colors at a position. On subsets of the prescribed
word $\mathtt{121}$, the maximal independent sets are $\{1,2\}$ and
$\{3\}$. The zero-rank-increment rule sends $\{3\}$ to all three
positions but sends $\{1,3\}$ to itself, so it is not a monotone closure.

\section[conditionalRankFlats]{Pinning and conditional flats}

Let $\tau$ be an independent partial assignment, put
$T_d=\{e:(e,d)\in\tau\}$, and let $U$ be the unpinned positions.
Write $M_c^\tau$ for the color-$c$ projection after pinning $\tau$
and deleting its positions. Its rank is $r_c-|T_c|$.

\begin{theorem}[Conditional ranks]
For every $S\subseteq U$,
\[
 \begin{gathered}
 \operatorname{rk}_{M_c^\tau}(S)\\
 =R(T_1,\dotsc,T_c\cup S,\dotsc,T_k)-|\tau|.
 \end{gathered}
\]
\end{theorem}

\begin{proof}
First pin only the other colors. These words maximize the score
$\sum_{d\ne c}|B(u,d)\cap T_d|$, with maximum $|\tau|-|T_c|$.
The slice theorem gives their color-$c$ projection rank after
subtracting this maximum from $R$.

In that projection, $T_c$ is independent because $\tau$ extends to
a word. Pinning color $c$ on $T_c$ contracts it, subtracting $|T_c|$
from the rank on $T_c\cup S$. The other pinned positions are loops
of the projection and can be removed.
\end{proof}

Thus the joint rank also determines all conditional flats. If $I$ is
independent in $M_c^\tau$, pinning color $c$ on $I$ gives
\[
 M_c^{\tau\cup(I\times\{c\})}=M_c^\tau/I.
\]
Its flats correspond to the flats of $M_c^\tau$ containing
$\operatorname{cl}_{M_c^\tau}(I)$: the usual contraction interval in
the lattice of flats. Mixed circuits become
\href{independence.htm#conditionalCircuits}{ordinary circuits after
the other colors are pinned}.

\section[branchQuotients]{Color branches and a common contraction}

With two colors, pinning a position gives ordinary deletion or
contraction in its color projection. With more colors, the deletion
side splits into several branches sharing one contraction.

Fix a position $e$ allowing color $c$ and at least one other color.
On $U=E\setminus\{e\}$, put
\[
 \begin{aligned}
 K&=M_c(\mathcal F/(e,c)),\\
 N_a&=M_c(\mathcal F/(e,a))\qquad(a\ne c),
 \end{aligned}
\]
using only nonempty \href{operations.htm#deletionAndContraction}{minors}.
The ranks are $r_c-1$ and $r_c$, respectively.

\begin{theorem}[Common contraction]
Let $D$ be a nonempty set of other colors possible at $e$. The union
\[
 \mathcal B(N_D)=\bigcup_{a\in D}\mathcal B(N_a)
\]
is the basis family of an ordinary matroid $N_D$. There is a matroid
$Q_D$ on $E$, with $e$ neither a loop nor a coloop, such that
\[
 Q_D\setminus e=N_D,\qquad Q_D/e=K.
\]
Thus $K$ is an \defin{elementary quotient} of every $N_D$.
\end{theorem}

\begin{proof}
Restrict the color at $e$ to $\{c\}\cup D$ and take the color-$c$
projection, calling it $Q_D$. Its bases avoiding $e$ are exactly the
displayed union. Its bases containing $e$, with $e$ removed, are
the bases of $K$. Both kinds occur by hypothesis.
\end{proof}

For every $S\subseteq U$, ordinary deletion and contraction give
\[
 \begin{gathered}
 \operatorname{rk}_{N_D}(S)=\max_{a\in D}\operatorname{rk}_{N_a}(S),\\
 \operatorname{rk}_K(S)\leq\operatorname{rk}_{N_D}(S)
 \leq\operatorname{rk}_K(S)+1.
 \end{gathered}
\]
Their closures satisfy $\operatorname{cl}_{N_D}(S)\subseteq
\operatorname{cl}_K(S)$, so
\[
 \mathcal L(K)\subseteq\mathcal L(N_D).
\]
This inclusion preserves intersections. The map
$X\mapsto\operatorname{cl}_K(X)$ is a surjection in the other
direction and preserves joins. These statements follow from the
closure formulas in $Q_D$: deletion uses $S$, while contraction uses
$S\cup\{e\}$. In particular, every $K$-flat is an $N_D$-flat, and
closing a union in $K$ computes the image of a join.

\begin{example}[The seven-word family]
Use the \href{independence.htm#mixedCircuits}{seven-word family} and
distinguish color $1$. Branch at position $1$, leaving
$U=\{2,3,4\}$. The remaining color-$1$ bases are
\[
 \begin{array}{c|c}
 u_1&\text{bases on }U\\ \hline
 1&\{2\},\{3\},\{4\}\\
 2&\{2,3\},\{3,4\}\\
 3&\{2,4\},\{3,4\}
 \end{array}
\]
The common contraction is $K=U_{1,3}$, while the union of the other
two basis families is $U_{2,3}$. In the color-$2$ branch, $3$ is a
coloop and $2,4$ are parallel; in the color-$3$ branch, $4$ is a
coloop and $2,3$ are parallel. Thus
\[
 \begin{aligned}
 \mathcal L(N_2)&=\{\varnothing,\{3\},\{2,4\},U\},\\
 \mathcal L(N_3)&=\{\varnothing,\{4\},\{2,3\},U\},\\
 \mathcal L(K)&=\{\varnothing,U\}.
 \end{aligned}
\]
The branch flats differ, while the flats of $K$ occur in both.
\end{example}

\section[valuatedCompletions]{Best completion scores}

Fix a color $c$ and real scores $a_{e,d}$ for $d\ne c$.
For a basis $B$ of $M_c$, define its \defin{best completion score}
\[
 \nu(B)\coloneqq
 \max_{\substack{u\in\mathcal F\\B(u,c)=B}}
 \sum_{e:u_e\ne c}a_{e,u_e}.
\]
Set $\nu(B)=-\infty$ on $r_c$-subsets with no completion.

\begin{theorem}[Valuated color projections]
The function $\nu$ is a \defin{valuated matroid} in the maximum
convention: for bases $B,C$ and $e\in B\setminus C$, some
$f\in C\setminus B$ satisfies
\[
 \begin{gathered}
 \nu(B)+\nu(C)\\
 \leq\nu(B-e+f)+\nu(C-f+e).
 \end{gathered}
\]
\end{theorem}

\begin{proof}
Choose best completions of $B,C$ and apply symmetric cycle exchange
through $e$. The cycle returns to color $c$ at a unique position
$f\in C\setminus B$. Its two resulting words have color-$c$ sets
$B-e+f$ and $C-f+e$. The total score is preserved, since each
position retains the same pair of entries across the two words.
Their scores are at most the best completion scores of their new
color classes, giving the inequality.
\end{proof}

Adding scores $b_e$ for color $c$, the bases maximizing
$\nu(B)+\sum_{e\in B}b_e$ form an ordinary matroid: project the
words maximizing the complete assignment score. These are the
\defin{initial matroids} of the valuation. For valuated-matroid
background, see D.~Speyer,
\href{https://arxiv.org/html/math/0410455v1}{\emph{Tropical Linear
Spaces}, Section 2}, using $-\nu$ to pass to the minimum convention.

For two colors, completions are unique and
$\nu(B)=\sum_{e\notin B}a_{e,d}$ for the other color $d$.
With more colors, the valuation can contain more information than
additive element weights. For example, take all words on four
positions with quotas $(2,1,1)$. Give score one to color $2$ at
positions $3,4$ and zero to all other assignments outside color $1$.
Then $\nu(\{3,4\})=0$ and every other color-$1$ basis has score one.
The three complementary pair sums are $1,2,2$, whereas additive
element weights would make them equal.

These best scores concern maxima over completions. The
\href{rank.htm#lorentzianConjecture}{completion-polynomial conjecture}
concerns sums over those completions and remains open.

\begin{question}[Rank axioms]
Which additional relations between color arguments characterize
the joint rank functions of colored matroids?
\end{question}
